\section{Related work}
\label{sec:related}

\subsection{Kernel generation as an evaluation substrate}

KernelBench~\cite{ouyang2025kernelbench} established the substrate this benchmark builds on:
a model is given a PyTorch module and asked for a faster implementation, and the submission is
graded by execution rather than by resemblance to a reference. KernelAscent inherits that task
and changes the question. KernelBench asks whether a model \emph{can} write an efficient kernel;
we ask whether that ability \emph{compounds} when the model trains on its own verified output.
The distinction matters for design rather than framing: a one-shot benchmark is free to use
tasks a strong model solves at parity, whereas a compounding benchmark is not, because a task
the control arm already solves contributes nothing to a contrast between arms
(Section~\ref{sec:instrument}).

Subsequent work has strengthened the grading contract along axes we depend on and do not
contribute to. KernelBench-Verified~\cite{kernelbenchverified} adds stronger input-distribution
checks and realistic precision settings; its authors are explicit that finite hidden tests are
not a proof of general correctness, and the same caveat applies to our grader.
SOL-ExecBench~\cite{solexecbench} evaluates against modelled hardware limits, which is the same
motivation as our per-task roofline ceiling~\cite{williams2009roofline}, though we use the
ceiling to normalise a score across heterogeneous tasks rather than to certify optimality.
FlashInfer-Bench~\cite{flashinferbench} and its trace schema~\cite{flashinfertrace} move in the
opposite direction from us, toward real inference workloads and runtime substitution of
generated kernels; that is the deployment question, and a benchmark that cannot yet show its own
measurement has range is not in a position to ask it.

At the system level, MLPerf Inference~\cite{mlperfinference} and vLLM's serving
benchmarks~\cite{vllmbench} measure end-to-end throughput and latency under declared quality
constraints. Our unit is different: we measure the change a self-improving agent produces in its
own kernel-writing ability, not the performance of a serving stack, and we make no claim of
MLPerf compliance.

\subsection{Self-improvement and self-modification}

The recursive self-improvement line this benchmark is aimed at spans several distinct mechanisms,
and the five rungs of Section~\ref{sec:benchmark} are an attempt to separate them rather than to
rank them. Voyager~\cite{voyager} accumulates an executable skill library, which is improvement
carried in \emph{stored artifacts} and is closest to our T3 procedure rung.
STOP~\cite{stop} recursively optimises the scaffolding around a code generator, and the Darwin
G\"odel Machine~\cite{dgm} maintains an archive of self-modifying variants; both improve the
\emph{procedure} while leaving weights fixed. HyperAgents~\cite{hyperagents} defines
improvement@k and a transfer experiment, which is the prior evaluation closest in spirit to the
contrasts reported here. Meta$^n$~\cite{metan} formulates recursion at the meta-agent level; we
cite it for positioning and do not independently verify its empirical claims.
AlphaEvolve~\cite{alphaevolve} is the existence proof that automated optimisation can reach
deployed computing systems, which is what makes the domain worth measuring carefully.

Our weight rung uses verified rejection-sampling fine-tuning, the method introduced by
STaR~\cite{zelikman2022star} and scaled by ReST~\cite{gulcehre2023rest}: sample, keep what
verifies, train on it, repeat. The contribution here is not the method but the controls around
it. A rising curve under this procedure is consistent with a better producer, with more data,
and with more sampling, and our three-arm decomposition exists to separate those.

The regularized objective in Section~\ref{sec:motivation:fed} adopts the KL-to-reference-policy
penalty standard in RLHF~\cite{ziegler2019klrlhf}, using the low-variance k3
estimator~\cite{schulman2020kl}. The novelty is the use rather than the term: our mechanism
analysis finds that sustained adapter drift and round-0 retention discriminate the runs that
improve, but finds it \emph{post hoc}, as properties of runs that already happened. A trust
region makes both of them settings.

\subsection{Why the instrument work is the contribution}

Two lines of work motivate the emphasis of this report. METR~\cite{metrrewardhacking} documents
frontier models manipulating tests and scoring mechanisms in automated research settings, which
is the case for grading a submission independently of the agent that produced it. We report a
failure that is adjacent but distinct, and in our view more common: nobody gamed anything here,
and the measurement still failed, because a scoring rule that cannot resolve the quantity under
study returns a confident null rather than an error. The infrastructure literature on
verification and isolation --- compute sanitizers~\cite{computesanitizer}, container
security~\cite{dockersecurity} --- addresses whether a harness can be trusted to run untrusted
code. Our finding is that a harness can be entirely trustworthy in that sense, execute every
candidate correctly, and still be incapable of measuring what the benchmark was built to
measure.

To our knowledge no prior self-improvement benchmark reports a dynamic-range precondition, a
throughput precondition, or a ledger of its own withdrawn results. We think all three should be
standard, and Section~\ref{sec:instrument} argues the case from our own failures rather than
from principle.
